% AY2013 材料力学 〔I〕（転記）
% 出典: 04_材料力学/original/04_材料力学/2013/2013za.pdf
% 要約: 片持ちはり（左端A固定）。中央B点（A から L/2）に集中荷重P,
%       右端C点（A から L）に集中荷重P と曲げモーメント M0(=5PL) が作用。
%       縦弾性係数E, 断面は幅b・高さhの長方形。支持反力, SFD/BMD,
%       B点上面のモール円・最大せん断応力, たわみ基礎式とC点たわみ角・たわみを問う。

\section*{〔I〕}

図に示すように片持ちはりの中央 B 点に集中荷重 $P$ が, 右端 C 点に
集中荷重 $P$ と曲げモーメント $M_0\ (=5PL)$ が負荷される問いに答えよ.
はりの縦弾性係数を $E$, 断面形状は幅 $b$, 高さ $h$ の長方形とする.

\begin{center}
\begin{tikzpicture}[>=Latex]
  % はり本体
  \draw[thick] (0,0) -- (6,0);
  % 固定壁（左端 A）
  \draw[thick] (0,-1.0) -- (0,1.0);
  \foreach \y in {-0.9,-0.6,...,0.95} {\draw (0,\y) -- (-0.28,\y+0.28);}
  % 節点
  \node at (0,-0.95) {A};
  \node at (3,-0.35) {B};
  \node at (6,-0.4) {C};
  % B点 集中荷重 P（下向き）
  \draw[->,very thick] (3,1.3) -- (3,0.05) node[midway,left]{$P$};
  % C点 集中荷重 P（下向き）
  \draw[->,very thick] (6,1.7) -- (6,0.05) node[above]{$P$};
  % C点 曲げモーメント M0（時計回り）
  \draw[->,thick] (6.35,0.55) arc (60:-60:0.55) node[midway,right]{$M_0$};
  % 寸法 L（A→C）
  \draw[<->] (0,1.45) -- (6,1.45) node[midway,fill=white]{$L$};
  % 寸法 L/2（A→B）
  \draw[<->] (0,-0.7) -- (3,-0.7) node[midway,fill=white]{$\dfrac{L}{2}$};
  % 座標 x
  \draw[->] (0,-1.1) -- (1.3,-1.1) node[right]{$x$};
\end{tikzpicture}
\end{center}

\begin{enumerate}[label=(\arabic*)]
  \item 支持点 A に生じる支持力および支持モーメントを図示し,
        つり合い式によりそれらを求めよ.

  \item せん断力 $Q(x)$ と曲げモーメント $M(x)$ を求め,
        せん断力図（SFD）と曲げモーメント図（BMD）を示せ.

  \item 中央 B 点での, はり上面におけるモールの応力円を図示し,
        最大せん断応力を求めよ.

  \item たわみ $v(x)$ に関するはりの基礎式を示し, 積分により
        たわみ角 $\dd v/\dd x=\theta(x)$ とたわみ $v(x)$ を求め,
        C 点におけるたわみ角およびたわみを求めよ.
\end{enumerate}
